\documentclass[a4paper,10pt]{article}
\usepackage[francais]{babel}
\usepackage{url}

\selectlanguage{francais}
\input{letter-icube-uds.sty}

\selectlanguage{francais}
\begin{document}


\name{Stéphane \textsc{Genaud}}
\fonction{Professeur des Universités}
\telephoneillkirch{+33 (0)3688 54542}
\email{genaud@unistra.fr}

\begin{letter}[Recommandation d'Etienne Michon pour l'offre d'ingénieur R\&D Cloud]
{à}
{Christine Morin, Yvon Jegou, Roberto Cascella}

Je connais  Etienne Michon depuis sa  deuxième année de master  en informatique,
pendant lequel  il a  décidé de faire  un stage recherche  avec Julien  Gossa et
moi-même sur un sujet concernant l'allocation  de ressources de clouds IaaS sous
forme de courtage pour les clients.  A  l'issue de son Master, nous avons obtenu
une allocation de recherche  de la DGA qui a permis à  Etienne de poursuivre son
travail sur le même sujet.

A environ  six mois de  la fin de  sa thèse, Etienne  a exploré le  périmètre de
recherche  convenu initialement.   Il s'agissait  de défricher  les possibilités
qu'une organisation  utilisatrice de  cloud peut
avoir pour  exécuter un ensemble  de tâches en  série sur une  infrastructure de
type IaaS tout en ayant un contrôle sur le
temps d'exécution et  le prix payé.  Pour cela, Etienne  est devenu le principal
contributeur    dans   le    développement   de    notre   logiciel    prototype
\textit{Schlouder},  qui  propose  différentes stratégies  d'ordonnancement  des
travaux.   Ce prototype  a  permis  de tester  l'exécution  de différents  types
d'application  (ensemble  de tâches  indépendantes  ou  workflows) et  plusieurs
plateformes. Les  tests préliminaires avaient  jusque là été réalisés  sur notre
infrastructure propre, et cette année, Etienne a pu étendre les expérimentations
à la plateforme européenne Bonfire.   L'ensemble des résultats expérimentaux qui
viennent   consolider    les   propositions   décrites   et    implantées   dans
\emph{Schlouder} sont actuellement rassemblés pour  la rédaction d'un article de
revue.

\vspace{.6cm}
Le travail  effectué par Etienne au  cours de sa thèse  est considérable.  Après
avoir conçu des  stratégies de courtage qui peuvent être  mises en {\oe}uvre, il
devient rapidement  nécessaire de  valider leurs  bénéfices.  Cela  nécessite un
effort  de  développement  important,  ainsi  que  de  nombreuses  manipulations
expérimentales qui ne peuvent être accomplies  que si le doctorant possède aussi
des  compétences avérées  d'ingénierie.  Nos  contraintes de  développements ont
amené  Etienne a  s'adapter  pour  réellement maîtriser  des  langages qu'il  ne
connaissait  que superficiellement  initialement (e.g  Perl). Il  s'est toujours
montré très investi et rigoureux dans ce domaine.  Grâce à ces expérimentations,
Etienne a pu démontrer la faisabilité de la solution logicielle proposée.

\vspace{1cm}
En conclusion, Etienne a mené son travail  de thèse avec sérieux sur un thème de
recherche  très  compétitif  du  fait  de  l'engouement  que  succite  le  cloud
computing. Son travail actuel doit logiquement déboucher sur la soutenance de la
thèse à l'automne. Je ne peux que  recommander vivement Etienne pour un poste de
R\&D dans  l'environnement attractif  proposé dans l'équipe-projet  Myriads.  Les
qualités  humaines  d'Etienne  font  qu'il, j'en  suis  persuadé,  s'intégrerait
facilement dans  un travail d'équipe  dans un domaine  dont il possède  une très
bonne connaissance. C'est pourquoi je recommande la candidature d'Etienne à cette
offre.

\end{letter}
\begin{flushright}
\includegraphics[width=.25\textwidth]{signgenaud.jpg}
\end{flushright}


\end{document}
